% Section 4: the shift from payroll to self-employment, seen in two federal
% series that count different things. The honesty constraint is not to present
% them as one workforce total.
\section{Payroll music jobs have fallen; solo creative businesses have risen}
\label{sec:payroll}

The Bureau of Labor Statistics has counted fewer payroll musician jobs in Texas
every survey vintage since 2005, and higher real hourly wages for the jobs that
remain.\footnote{The Occupational Employment and Wage Statistics program
(OEWS) is a semi-annual establishment survey the U.S.\ Bureau of Labor
Statistics runs jointly with each state workforce agency. Sampled employers
report the occupation and wage of every worker on their payroll; a
self-employed musician, a musician paid off the books, and a musician
performing under an independent-contractor arrangement never enter the
sampling frame at all. This section therefore reads OEWS as a payroll-only
ceiling on measured wage-and-salary music work rather than as a count of the
workforce. Program documentation at \url{https://www.bls.gov/oes/oes_ques.htm};
the tables read here are the May 2005 through May 2025 state and metropolitan
estimates for Musicians and Singers, Standard Occupational Classification code
27-2042, published at \url{https://www.bls.gov/oes/tables.htm}. BLS does not
publish an annual wage for this occupation in most cells because the work is
too irregular to annualise from an hourly rate.} On its own that pattern reads
as an occupation getting smaller and better paid. Set against the household
microdata in Section~\ref{sec:earnings}, a second reading opens: the same
work has moved out of payroll and into self-employment, which no payroll
wage series covers. The last subsection here tests that reading against the
timing of the two series and finds it weaker than a first look suggests.

In Texas, OEWS records \NEmpPayrollMusiciansTxFirst\ payroll musician jobs in
2005 and \NEmpPayrollMusiciansTxLast\ in 2025, a change of
\NEmpPayrollPctchgMusiciansTx\ over twenty years. The same program records a
national change of \NEmpPayrollPctchgMusiciansUs\ over the same window. For
the payroll jobs that remained in Texas, the real median hourly wage rose
\NWageHourlyPctchgMusiciansTx, from \NWageHourlyRealMusiciansTxFirst\ in 2005
to \NWageHourlyRealMusiciansTxLast\ in 2025, both in 2025 dollars
(Figure~\ref{fig:wagetrend}).\footnote{Every percentage change quoted in
this section is an endpoint change over the whole span,
\(100\,(X_T / X_0 - 1)\), where \(X_0\) is the first year with a published
value and \(X_T\) the last, so it is not an annual rate and it says nothing
about the path between the two years. Job counts are compared as published,
because a count of jobs is not a dollar figure and nothing in it needs
deflating. The wage is deflated before it is differenced,
\(w^{\mathrm{real}}_y = w_y \times d_y\) with \(d_y = C_{2025}/C_y\) the ratio
of national consumer price index annual averages, so the change reported here
is the change in purchasing power and not the change in the posted wage. Where
a series has no published value for a year, the endpoints move inward to the
first and last years that do, which for the Austin musician wage shortens the
window. These are the OEWS program's own published medians and counts; BLS
publishes no margin of error for these cells, so none is shown.}

\begin{figure}[htbp]
\centering
\includegraphics[width=\textwidth]{fig07_real_wage_trend.png}
\caption{Real median hourly wage for selected occupations, Austin metro and
Texas, 2005--2025, in 2025 dollars. Source: U.S.\ Bureau of Labor Statistics,
Occupational Employment and Wage Statistics, May 2005 through May 2025 tables
for Musicians and Singers (SOC 27-2042), deflated with the national consumer
price index (CPI-U, annual average). How to read it: the axis does not start
at zero. The Austin musician line swings widely because it rests on between
about 100 and 700 payroll jobs depending on the year, so movements of several
dollars between years reflect sampling variation rather than changes in what
Austin musicians are paid. Caveat: OEWS excludes the self-employed by design,
so this series describes a small and better-paid slice of the occupation
rather than the occupation itself.}
\label{fig:wagetrend}
\end{figure}

\subsection{Solo creative businesses grew while payroll music jobs fell}

The U.S.\ Census Bureau publishes a separate account of solo creative work
through Nonemployer Statistics, an administrative tabulation the Bureau builds
from federal tax filings rather than from a survey.\footnote{Nonemployer
Statistics counts business establishments with no paid employees and at least
one thousand dollars of annual receipts (one dollar in construction), drawn
from the Internal Revenue Service's business tax records: Form 1040 Schedule C
for sole proprietors, Form 1065 for partnerships, and the corporate returns
1120 and 1120S. The sampling frame is therefore tax filings, not households
or payrolls, so the file captures registered solo businesses that OEWS
structurally cannot see but omits cash gig work that never enters the tax
record. The industry code used here, NAICS 7115, is Independent Artists,
Writers, and Performers, a category broader than musicians alone. County-level
release: \url{https://www2.census.gov/programs-surveys/nonemployer-statistics/datasets/2023/historical-datasets/nonemp23co.zip}.
Program documentation:
\url{https://www.census.gov/programs-surveys/nonemployer-statistics.html}.
The Census Bureau protects county cells with noise infusion, adding small
random errors to published values instead of suppressing them, which matters
more for a county cell than for a state one.} That series counts a different
unit (a solo business, not a job) over a different universe (all independent
artists, writers and performers, not musicians alone) than OEWS. The two
files are therefore not additive, and Figure~\ref{fig:substitution} draws
them as separate panels so no reader adds them.

\begin{figure}[htbp]
\centering
\includegraphics[width=\textwidth]{fig06_payroll_vs_nonemployer.png}
\caption{Payroll music jobs against independent-artist businesses without
employees, each indexed to its own first available year. Sources: U.S.\
Bureau of Labor Statistics, OEWS, for payroll employment in Musicians and
Singers (SOC 27-2042); U.S.\ Census Bureau, Nonemployer Statistics, for
establishments in NAICS 7115 (Independent Artists, Writers, and Performers).
How to read it: the two series count different things over different
universes, jobs against businesses and musicians against all independent
artists, so the figure shows the direction each is moving and nothing more.
The panels must not be added or netted.}
\label{fig:substitution}
\end{figure}

Nonemployer establishments in NAICS 7115 rose
\NSubstitutionNonemployerTravisPctchg\ in Travis County between 2012 and
2023 and \NSubstitutionNonemployerTxPctchg\ across Texas over the same
window.\footnote{Nonemployer Statistics reports Travis County
establishments in NAICS 7115 at 5{,}382 in 2012 and 8{,}696 in 2023, and
Texas statewide at 45{,}021 and 74{,}834 over the same years. Both figures
are business counts, not people, and the Bureau's noise-infusion procedure
means a single county cell has a slightly larger measurement error than the
state total.} Texas payroll music jobs, measured over their own longer OEWS
window of 2005 to 2025, changed
\NEmpPayrollPctchgMusiciansTx.\footnote{Figure~\ref{fig:substitution}
rebases each series on its own first available year,
\(I_{s,t} = 100 \times X_{s,t} / X_{s,t_0(s)}\), where \(X_{s,t}\) is the
level of series \(s\) in year \(t\) and \(t_0(s)\) is that series' first
year: 2005 for OEWS payroll jobs, 2012 for the Nonemployer Statistics
establishment count. Rebasing is what allows two series measured in
different units, jobs in one case and businesses in the other, to sit on a
common vertical scale, and it is also the reason the two panels cannot be
added or netted: an index number carries no unit and no count, only a
direction and a rate of movement relative to its own starting point.} These
two series do not count the same people. This study makes no estimate of
the total size of the Austin music workforce, because no public source
supports one.

Among the four large Texas counties compared here, Travis County has the
highest rate of independent-artist businesses relative to its population:
\NNesEstabPerOneZerokTravisTwoZeroTwoThree\ nonemployer establishments in
NAICS 7115 per 10{,}000 residents in 2023, against
\NNesEstabPerOneZerokDallasTwoZeroTwoThree\ in Dallas,
\NNesEstabPerOneZerokBexarTwoZeroTwoThree\ in Bexar and
\NNesEstabPerOneZerokHarrisTwoZeroTwoThree\ in Harris
(Figure~\ref{fig:metrones}).\footnote{The rate is
\(10^{4} \times E_{c,y} / N_{c,y}\), where \(E_{c,y}\) is the count of
Nonemployer Statistics establishments in NAICS 7115 in county \(c\) and
year \(y\), and \(N_{c,y}\) is that county's resident population in the
same year, taken from the U.S.\ Census Bureau's Vintage 2025 Population
Estimates Program release
(\url{https://www2.census.gov/programs-surveys/popest/datasets/2020-2025/counties/totals/}).
Both sides share a year, so the measure is a rate rather than a count and
counties of different sizes can sit beside one another. The numerator
counts businesses and not people, and one person can hold more than one,
so this rate is a proxy for how common solo creative work is rather than
a headcount of solo creatives. Four of Texas's 254 counties appear here,
so the ranking is bounded by those four.}

\begin{figure}[htbp]
\centering
\includegraphics[width=0.86\textwidth]{fig08_metro_nonemployer.png}
\caption{Independent-artist businesses without employees per 10,000
residents, four Texas counties, 2020--2023. Sources: U.S.\ Census Bureau,
Nonemployer Statistics, NAICS 7115; county population from the U.S.\ Census
Bureau's Vintage 2025 Population Estimates Program. How to read it: a
higher line means more solo creative businesses relative to population.
Caveats: NAICS 7115 covers independent artists, writers and performers
together, so the count is a proxy for solo creative work rather than a
headcount of musicians; the series is shorter than the 2012--2023
establishment counts because the population denominator comes from the
Vintage 2025 estimates, which do not reach back to 2012; and four of
Texas's 254 counties appear here, so the ranking is bounded by those
four.}
\label{fig:metrones}
\end{figure}

Travis County's concentration of solo creative businesses is real; what
those businesses take in is less clear. Nonemployer Statistics reports
gross receipts before any expense: real receipts per Travis County
establishment in NAICS 7115 rose \NNesReceiptsPerEstabPctchgTravis\
between 2012 and 2023, to
\NNesReceiptsPerEstabTravisLast\ in 2025
dollars.\footnote{Receipts per business is
\(\bar{r}_{c,y} = R^{\mathrm{real}}_{c,y} / E_{c,y}\), where
\(R^{\mathrm{real}}_{c,y}\) is total NAICS 7115 receipts in county \(c\)
and year \(y\), deflated to 2025 dollars by \(d_y = C_{2025}/C_y\), and
\(E_{c,y}\) is the establishment count. This is an arithmetic mean with no
distribution behind it: the Census Bureau publishes only totals and counts
for this category, so a small number of high-earning businesses can carry
the average and no median is available. Noise infusion further complicates
year-to-year comparisons at the county level.} Alongside that receipts
figure sits a separate reading of the cost side of solo music work. The
2022 Greater Austin Music Census, a self-selected online survey of Austin
music workers with no probability sample behind it, asked respondents to
report annual spending on their own music business across eight categories;
the \NCensusSpendCompleteN\ music creatives who answered every one of those
eight categories reported a median of \NCensusSpendTotalMedian\ a year in
2025 dollars.\footnote{Sound Music Cities and City of Austin Music \&
Entertainment Division, \textit{Greater Austin Music Census}, 2022;
respondent-level microdata downloaded from City of Austin open data,
dataset \texttt{an3p-3yqx}, at
\url{https://data.austintexas.gov/}, retrieved 1 August 2026. The
instrument was distributed through music-community mailing lists and city
channels between 15 July and 12 September 2022, drew 2{,}227 respondent
records, and carried no probability weights. The complete-spender base of
\NCensusSpendCompleteN\ is 22.5 percent of respondents; the eight
categories are gear and rentals, publicity and promotion, rehearsal or
work space, new recordings, merchandise, supplies, web and social media,
and accounting and legal. The median above is the median of respondent
totals, not a sum of category medians.} The receipts figure is Census
Bureau administrative data on all Travis County solo businesses in NAICS
7115; the spending figure is a self-selected survey of Austin music
creatives specifically. The two describe different populations in
different years and cannot be subtracted from one another. They sit
together here because a receipts figure alone overstates what a solo
creative keeps.

\subsection{The substitution reading is an inference, not a tracked transition}

The two series do not move in the same years, which weakens the reading
that payroll music jobs became solo creative businesses rather than
supporting it. Most of the Texas OEWS payroll decline happened between 2005
and 2012, and the payroll count has moved sideways since; the Nonemployer
Statistics series begins in 2012, after that fall. What the two files show
together is a payroll count far below its 2005 level and a nonemployer
establishment count that has risen since 2012. Alongside them sits the ACS
PUMS microdata finding that \NPumsAustinMusSelfemp\
(\NPumsAustinMusSelfempMoe) of Austin musicians worked for themselves over
2020 to 2024, which is a single pooled cross-section from a household
survey and carries no time path of its own.\footnote{Musicians and Singers
(OCCP 2752), Austin-Round Rock-San Marcos metro, ages 18--64, employed
(ESR 1/2/4/5), classified self-employed if Class of Worker equals 6
(unincorporated) or 7 (incorporated), from the U.S.\ Census Bureau's
American Community Survey 2020--2024 five-year Public Use Microdata
Sample for Texas
(\url{https://www2.census.gov/programs-surveys/acs/data/pums/2024/5-Year/csv_ptx.zip}).
The share rests on 77 unweighted records; the margin of error uses the
Bureau's successive-difference replication method against the 80
replicate weights.} No public dataset follows individual musicians across
employment arrangements, so nothing here shows that particular OEWS
payroll jobs became particular Nonemployer Statistics businesses.

What the evidence does establish is narrower and still consequential: the
number of music jobs any payroll statistic can observe is lower now than in
2005, while the count of solo creative businesses in the same broad
category has risen since 2012, so payroll records describe less of Austin's
music work than they once did. No public source estimates what fraction of
Austin's music workforce now sits outside payroll records altogether,
because no source combining the payroll and self-employed sides at the
metro level exists. Section~\ref{sec:venues} turns from workers to
workplaces, and tracks the Austin live-music rooms where those solo careers
are performed, using the state Comptroller's monthly alcohol-tax filings as
a room-level receipts record the payroll surveys cannot supply.
